\section{Measure the actual data-access sequence}
Use a cold cache at function entry with sixteen lines, sixteen bytes per line,
and least-recently-used replacement. Compare fully associative, direct-mapped and
two-way placement at the same 256-byte capacity. Stop at the supplied return
continuation, before the caller stores \texttt{v0} in \texttt{completed}. Include
the function's descriptor loads and frame traffic; exclude assembly, session
creation, the caller and debugger inspection.

A completed simulated instruction and a data-access notice are different units.
In the current core, an aligned \texttt{ld} performs two four-byte reads and an
aligned \texttt{sd} performs two four-byte writes. Each notice is counted separately
by the Data Cache Simulator. The first notice can miss and the second hit in the
same line. This is an implementation detail of the observation stream, not two
architectural load instructions.

Supply an eight-byte-aligned stack for these comparisons. The current core handles
an unaligned doubleword through byte accesses, changing the notice count and address
sequence. Even though both paths can produce the same saved register values, mixing
stack alignments would make the table compare more than tile size. The measurement
fixture sets \texttt{sp} to \texttt{0x7fff0000} for every call.

The model places each notice using its starting address. It does not split a
notice that crosses a cache-line boundary. Here all input word accesses and both
word halves of aligned output/frame doublewords fit within sixteen-byte lines.
These choices make the address interpretation explicit. The model has no memory
latency, prefetcher or simulated pipeline and supplies no cycle estimate.

\subsection{Predict the baseline's notice count}
For nonempty output, the baseline makes \(2MNK\) input word reads and \(2MN\)
output word writes. Saving and restoring nine full registers contributes thirty-six
word notices; loading the descriptor contributes four. Thus its whole-function
count is
\[
  2MNK+2MN+40.
\]
For \(M=K=N=16\), this is 8744. An empty output still has the forty frame and
descriptor notices in this implementation. These counts describe the illustrated
program. An equivalent solution that uses different private storage can have a
different whole-function count.

The blocked version still reads two input words per product. For nonzero \(K\),
it additionally zeroes output, then reads and writes each output doubleword once
per reduction tile. Its notice count is
\[
  2MNK+2MN+4MN\left\lceil K/T\right\rceil+40.
\]
Fewer cache misses therefore need not imply fewer accesses. Output partial sums
add traffic while the changed ordering can preserve useful input lines.
